%! TEX root = main.tex

\section{Numerical Dispersion Analysis}
\label{sc:14}

It is well known that spatial discretization introduces a dispersion error: the discrete solution does not propagate with the same phase velocity as the continuous (physical) wave. In the frequency domain, this phenomenon is conveniently quantified by comparing the continuous wavenumber \(k\) with the discrete (numerical) wavenumber \(k_h\), associated with the finite-dimensional approximation. The resulting discrepancy represents the numerical dispersion of the method.

\subsection{Eigenvalue Problem Setting}

To isolate dispersion effects from external excitations, the homogeneous problem is considered by setting \(f = 0\) and imposing homogeneous boundary data. In particular, a homogeneous Neumann condition is prescribed at the left boundary, namely \(g_N = 0\), which implies \(\mu u'(0) = 0\), while a homogeneous Dirichlet condition is imposed at the right boundary, \(u(L) = 0\). Under these assumptions, the Helmholtz model reduces to an eigenvalue problem for the spatial operator.

\subsection{Continuous and Discrete Formulations}

In the continuous setting and for constant coefficients \(\rho = \mu = 1\), the eigenmodes satisfy the classical second-order equation

\begin{equation}
    -u''(x) = k^2 u(x)
    \label{eq:cont_eig}
\end{equation}

where \(k^2\) denotes the exact eigenvalues of the continuous problem and \(k\) is the corresponding continuous wavenumber. By applying the finite element discretization derived in section \ref{sc:11} and restricting the solution to the discrete space, the homogeneous algebraic system assumes the form of a generalized eigenvalue problem. Denoting by \(A\) the stiffness matrix and by \(M\) the mass matrix, the discrete modes satisfy

\begin{equation}
    \mathbf{A}\underline{u} = k_h^2 \mathbf{M}\underline{u}
    \label{eq:disc_eig}
\end{equation}

where \(k_h^2\) are the discrete eigenvalues and \(k_h\) is the associated numerical wavenumber. The dispersion error is therefore interpreted as the difference between the spectra of the continuous and discrete operators, expressed in terms of wavenumbers.

\subsection{Numerical Results}

The accuracy of the method is assessed by introducing the dimensionless resolution parameter \(h/\lambda\), where \(h\) denotes the characteristic mesh size and \(\lambda\) the wavelength. The numerical dispersion is represented through the dispersion curve, defined as the ratio \(k_h/k\) plotted as a function of \(h/\lambda\). The generalized eigenvalue problem \eqref{eq:disc_eig} is solved in MATLAB for a sequence of meshes with \(N \in \{10,\ 20,\ 40,\ 80,\ 160\}\) elements. For each discrete eigenvalue \(k_h^2\), the corresponding exact wavenumber \(k\) is identified analytically and the ratio \(k_h/k\) is computed. The resulting dispersion curve is reported in figure \ref{fig:dispersion}, together with the asymptotic Taylor expansion
\begin{equation}
    \frac{k_h}{k} \approx 1 + \frac{\pi^2}{6}\left(\frac{h}{\lambda}\right)^2
    \label{eq:taylor_disp}
\end{equation}
which holds for the consistent-mass \(\mathbb{P}_1\) FEM in the well-resolved regime \(h/\lambda \ll 1\).

\begin{figure}[!htb]
    \centering
    \includegraphics[width=0.90\textwidth]{00_Images/HW1/4_Dispersion.png}
    \caption{Numerical dispersion curve \(k_h/k\) as a function of the resolution parameter \(h/\lambda\) for the consistent-mass \(\mathbb{P}_1\) FEM (blue dots). The dashed black line at \(k_h/k = 1\) is the exact non-dispersive reference, while the solid parabola represents the asymptotic Taylor expansion \(1 + \frac{\pi^2}{6}(h/\lambda)^2\). For well-resolved meshes the two curves coincide, whereas near the Nyquist limit the exact discrete symbol departs from the parabola, reaches a local maximum and settles at approximately \(k_h/k \approx 1.10\) at \(h/\lambda = 0.5\)}
    \label{fig:dispersion}
\end{figure}

\subsection{Physical Interpretation}

The dispersion curve in figure \ref{fig:dispersion} reveals the characteristic overestimation behavior of the consistent-mass \(\mathbb{P}_1\) FEM: unlike lumped-mass formulations, which underestimate the physical wavenumber, the consistent-mass scheme yields \(k_h > k\), so that the ratio \(k_h/k\) is strictly greater than unity throughout the resolved spectrum. Three distinct resolution regimes can be identified.

\paragraph{Fine meshes: \(h/\lambda \to 0\)} In the well-resolved regime, the numerical results follow the asymptotic parabola \eqref{eq:taylor_disp} with remarkable accuracy, confirming the theoretical second-order convergence of the dispersion error. The overestimation of the wavenumber is practically negligible: below \(h/\lambda \approx 0.05\), corresponding to 20 or more elements per wavelength, the deviation from unity is smaller than \(0.1\%\) and the numerical phase velocity coincides with the physical one to engineering accuracy.

\paragraph{Intermediate resolution} At \(h/\lambda = 0.1\), i.e.\ approximately 10 elements per wavelength, the ratio \(k_h/k \approx 1.017\), corresponding to a dispersion error of about \(1.7\%\). This level of resolution is commonly considered a practical minimum for engineering applications where moderate accuracy is acceptable. The asymptotic expansion \eqref{eq:taylor_disp} remains in close agreement with the exact discrete symbol up to this range.

\paragraph{Coarse meshes: approach to the Nyquist limit} As \(h/\lambda\) increases beyond \(0.2\), the exact dispersion curve progressively departs from the parabolic approximation, which diverges monotonically, whereas the physical discrete symbol reaches a local maximum and subsequently decreases. Near the Nyquist limit \(h/\lambda \to 0.5\), the curve settles at a finite value of approximately \(k_h/k \approx 1.10\). This behavior, captured in full by the numerical computation since the MATLAB solver resolves the complete discrete spectrum without restriction, demonstrates that the parabolic expansion is a low-frequency asymptotic and does not bound the solution at high wavenumbers. Points in the plot that momentarily exceed the upper axis limit before returning to \(k_h/k \approx 1.10\) are therefore physically meaningful and not artifacts: they correspond to the exact symbol of the discrete operator evaluated in the poorly resolved regime.

\paragraph{General rule of thumb} The analysis confirms that at least 10 elements per wavelength (\(h/\lambda \leq 0.1\)) are required to keep the dispersion error below \(2\%\) and that 20 or more elements per wavelength (\(h/\lambda \leq 0.05\)) are needed for high-fidelity simulations. These thresholds are consistent with the standard engineering guidelines for \(\mathbb{P}_1\) FEM applied to wave propagation problems.